% =====================================================================
%  FICHE 11 — THEME 3 — Corrigé MOYEN — Équilibrer une équation redox
% =====================================================================
\documentclass[11pt,a4paper]{article}
\usepackage[utf8]{inputenc}
\usepackage[T1]{fontenc}
\usepackage{lmodern}
\usepackage[french]{babel}
\usepackage{amsmath,amssymb}
\usepackage[version=4]{mhchem}
\usepackage{siunitx}
\sisetup{output-decimal-marker={,},per-mode=symbol}
\usepackage{xcolor}
\usepackage{array}
\usepackage{booktabs}
\usepackage{enumitem}
\usepackage[most]{tcolorbox}
\usepackage[a4paper,margin=2.0cm,headheight=26pt,headsep=10pt]{geometry}
\usepackage{fancyhdr}
\usepackage[hidelinks]{hyperref}

\definecolor{primary}{RGB}{146,95,160}
\definecolor{primarydark}{RGB}{92,52,104}
\definecolor{corrcol}{RGB}{0,135,90}
\newcommand{\NO}{\ensuremath{\mathrm{N.O.}}}

\newcounter{exo}
\newtcolorbox{corr}[1][]{enhanced,breakable,boxrule=0.9pt,arc=2pt,colback=corrcol!4,
  colframe=corrcol,fonttitle=\bfseries,coltitle=white,left=3mm,right=3mm,top=2.2mm,bottom=2.2mm,
  title={Corrigé~\refstepcounter{exo}\theexo\ifx\relax#1\relax\relax\else\ \textmd{--- \itshape #1}\fi}}

\pagestyle{fancy}\fancyhf{}
\fancyhead[L]{\small\textcolor{primarydark}{\textbf{Fiche 11 · Thème 3} — Corrigé}}
\fancyhead[R]{\small\textcolor{primarydark}{\textsc{Moyen}}}
\fancyfoot[C]{\small\thepage}
\renewcommand{\headrulewidth}{0.8pt}
\renewcommand{\headrule}{\hbox to\headwidth{\color{primary}\leaders\hrule height \headrulewidth\hfill}}

\begin{document}
\begin{center}
{\LARGE\bfseries\color{primarydark}Thème 3 — Corrigé Moyen}\\[-2pt]
{\color{corrcol}\rule{\textwidth}{1pt}}
\end{center}

\begin{corr}[assembler deux demi-équations — milieu acide]
La première échange 3 e$^-$, la seconde 1 e$^-$ : on multiplie la seconde par 3, puis on additionne.
Les 3 électrons s'éliminent :
\[
\boxed{\ce{NO3^- + 3 Fe^{2+} + 4 H+ -> NO + 3 Fe^{3+} + 2 H2O}}
\]
\textbf{Contrôle des charges :} gauche $-1 + 3(+2) + 4(+1) = +9$ ; droite $0 + 3(+3) + 0 = +9$. \checkmark
Atomes : N 1/1, Fe 3/3, O 3/3, H 4/4. \checkmark
\end{corr}

\begin{corr}[équilibrer avec des ions hydronium]
\begin{enumerate}[label=(\alph*),leftmargin=2em,topsep=1pt,itemsep=1.5pt]
  \item Couples : \ce{PbO2}/\ce{Pb^{2+}} (Pb : $+\mathrm{IV}\to+\mathrm{II}$, réduction, $n=2$) et
        \ce{Cr2O7^{2-}}/\ce{Cr^{3+}} (Cr : $+\mathrm{III}\to+\mathrm{VI}$, oxydation, $2\times3=6$ e$^-$).
  \item Demi-équations (milieu acide) :
        \[
        \ce{PbO2 + 4 H+ + 2 e^- -> Pb^{2+} + 2 H2O} \quad ; \quad
        \ce{2 Cr^{3+} + 7 H2O -> Cr2O7^{2-} + 14 H+ + 6 e^-}.
        \]
        On multiplie la réduction par 3 (6 e$^-$) et on additionne :
        \[
        \ce{3 PbO2 + 2 Cr^{3+} + H2O -> 3 Pb^{2+} + Cr2O7^{2-} + 2 H+}.
        \]
        (après simplification : $12$ vs $14$ \ce{H+} $\to 2$ \ce{H+} à droite ; $7$ vs $6$ \ce{H2O} $\to 1$ à gauche).
  \item On remplace les 2 \ce{H+} (droite) par 2 \ce{H3O+} en ajoutant 2 \ce{H2O} à gauche :
        \[
        \boxed{\ce{3 PbO2 + 2 Cr^{3+} + 3 H2O -> 3 Pb^{2+} + Cr2O7^{2-} + 2 H3O+}}
        \]
        Le \textbf{coefficient de \ce{H3O+} est 2}. \emph{(Réponse C, annales 2024-Q7.)}
        Contrôle : charges $2(+3)=+6$ à gauche $=3(+2)+(-2)+2(+1)=+6$ à droite. \checkmark
\end{enumerate}
\end{corr}

\begin{corr}[une demi-équation en milieu basique]
Mn : $+\mathrm{VII}\to+\mathrm{IV}$, $n=3$. En \textbf{acide} :
$\ce{MnO4^- + 4 H+ + 3 e^- -> MnO2 + 2 H2O}$. On ajoute 4 \ce{OH^-} des deux côtés ; les
$4(\ce{H+}+\ce{OH^-})$ deviennent 4 \ce{H2O}, on simplifie 2 \ce{H2O} :
\[
\boxed{\ce{MnO4^- + 2 H2O + 3 e^- -> MnO2 + 4 OH^-}}
\]
Contrôle : charges gauche $-1-3=-4$, droite $4(-1)=-4$ \checkmark ; O : $4+2=6=2+4$ ; H : $4=4$. \checkmark
\end{corr}

\begin{corr}[permanganate et eau oxygénée]
Couples : \ce{MnO4^-}/\ce{Mn^{2+}} ($n=5$) et \ce{O2}/\ce{H2O2} (O : $-\mathrm{I}\to 0$, oxydation).
\[
\ce{MnO4^- + 8 H+ + 5 e^- -> Mn^{2+} + 4 H2O} \quad ; \quad \ce{H2O2 -> O2 + 2 H+ + 2 e^-}.
\]
On multiplie la réduction par 2 (10 e$^-$) et l'oxydation par 5 (10 e$^-$), puis on additionne et on
simplifie les \ce{H+} ($16-10=6$) :
\[
\boxed{\ce{2 MnO4^- + 6 H+ + 5 H2O2 -> 2 Mn^{2+} + 8 H2O + 5 O2}}
\]
Coefficients : $2,6,5 \to 2,8,5$ (\textbf{réponse B}, annales 2022-Q18). Contrôle : O $8+10=18=8+10$ ;
H $6+10=16=16$ ; charges $2(-1)+6=+4=2(+2)$. \checkmark
\end{corr}

\begin{corr}[somme des coefficients]
\begin{enumerate}[label=(\alph*),leftmargin=2em,topsep=1pt,itemsep=1.5pt]
  \item Couples : \ce{Sn}/\ce{SnO2} (Sn : $0\to+\mathrm{IV}$, oxydation, $n=4$) et \ce{NO3^-}/\ce{NO2}
        (N : $+\mathrm{V}\to+\mathrm{IV}$, réduction, $n=1$).
        \[
        \ce{Sn + 2 H2O -> SnO2 + 4 H+ + 4 e^-} \quad ; \quad \ce{NO3^- + 2 H+ + e^- -> NO2 + H2O}.
        \]
        On multiplie la réduction par 4 et on additionne :
        \[
        \boxed{\ce{Sn + 4 NO3^- + 4 H+ -> SnO2 + 4 NO2 + 2 H2O}}
        \]
  \item Somme des coefficients des \textbf{réactifs} : $1+4+4=\mathbf{9}$.
        Somme des \textbf{produits} : $1+4+2=\mathbf{7}$.
        La somme des réactifs ($9$) est \textbf{supérieure} à celle des produits ($7$).
        \emph{(Réponse A, livre QCM~13.)}
\end{enumerate}
\end{corr}

\end{document}
